\chapter{机器学习与 AI Alpha 分层答案}

\begin{enumerate}[leftmargin=*]
  \item 对训练样本 $D_{\rm train}$，经验风险是 $\widehat R(f)=n^{-1}\sum_i\ell(f(x_i),y_i)$；泛化风险是未来分布下的 $R(f)=\mathbb E_{(X,Y)\sim\mathcal P_{\rm future}}\ell(f(X),Y)$；正则化目标是 $\widehat R(f)+\lambda\Omega(f)$。三者的积分域和优化对象不同，不能只用训练损失回答未来表现。
  \item 预测损失只比较数值误差，校准检查 $\mathbb P(Y=1\mid p\approx u)\approx u$，解释稳定性检查数据/时间扰动下重要性是否保持。一个模型可以低 MSE 但概率过度自信，也可以校准良好但交易信号弱；因果结论还需要处理定义、识别和反事实设计。
  \item 先按题目给的四个 $(y_i,\hat y_i)$ 计算残差平方并平均，得到训练风险 $0.4256$；把同一模型放到冻结推理段得 $1.5154$；在训练目标加 $\lambda\Omega=0.0100$ 得 $0.4356$。逐项保留分子、样本数和惩罚项，避免把测试误差误叫训练误差。
  \item Brier 分数是 $n^{-1}\sum_i(p_i-y_i)^2$。校准前逐项平方误差相加再除以样本数得 $0.16$，校准映射后同样步骤得 $0.025$；改善只说明概率刻度更接近标签，不说明排序、换手或成本后收益也改善。
  \item 四步窗口的输入 shape 从 $(B,3,F)$ 变为 $(B,4,F)$；padding mask 需从 $(B,3)$ 同步变为 $(B,4)$，并把目标写成严格晚于第四个观测的 $y_{t+4}$。先打印每个 batch 的最后有效索引，再断言标签时间大于窗口末端，才能排除错位而非只看 shape。
  \item 独立 ledger 先记录输入列顺序、窗口、目标和 expected，再运行模型生成 observed；修改第三笔交易或 schema 后只更新临时 fixture，比较哈希和输出。若 expected 由同一实现循环生成，任何共同 bug 都会被掩盖，不能算验证。
  \item 随机切分使同一资产的相邻日期和重叠标签跨越训练/测试；全样本 scaler 的均值方差含未来；推理段 early stopping 又把测试反馈变成选择信号。修复是按决策时间拟合预处理、purge 标签重叠、冻结外层测试。
  \item 先冻结漂移阈值和响应等级：轻微输入校准偏移可在训练窗内重校准；概念漂移需停用并进入重训审批；执行数据异常则降仓或停机。每次动作保留触发统计、窗口、版本、批准人和恢复条件，不能事后按收益选择动作。
\end{enumerate}
